\section{Related work}
\subsection{Diffusion-Based Models in Image Synthesis.}
Diffusion models achieve high-fidelity image synthesis through iterative Markov chain diffusion steps. Latent Diffusion Models (LDMs)~\cite{rombach2022high} shift this process from pixel space to compressed latent space, reducing computational overhead while maintaining image quality, fostering widespread adoption. Dhariwal~\cite{dhariwal2021diffusion} further demonstrated that diffusion models can be conditioned on text, facilitating the generation of semantically coherent and contextually relevant images. Imagic and its variants~\cite{kawar2023i, brack2024ledits++} improve control through joint optimization of images and text descriptions, enabling semantic-level manipulation. Recent findings~\cite{lian2024llmgrounded, wu2025paragraph} indicate that longer, more descriptive captions enhance generative quality. Motivated by this, we re-caption the nuScenes dataset with detailed scene descriptions to improve conditional image synthesis.

In addition, numerous diffusion-based frameworks are proposed for conditional generation with explicit spatial control through external signals.
ControlNet~\cite{zhang2023adding} augments pretrained diffusion backbones with auxiliary neural networks that integrate layout cues, segmentation maps, or depth priors to achieve spatially aligned synthesis.
GLIGEN~\cite{li2023gligen} introduces a gated self-attention mechanism to ground image generation on bounding boxes while maintaining consistency with textual prompts.
In contrast, T2I-Adapter~\cite{mou2024t2i} employs lightweight, modular adapters that inject conditioning signals into frozen diffusion models via element-wise operations.
Its simplicity and efficiency make T2I-Adapter particularly suitable for our decoupled learning scheme based on geometric layout conditions.
\subsection{Controllable Editing Driving Scenes.}
The growing demand for large-scale annotated driving-scene datasets, coupled with high manual labelling costs, necessitates efficient synthetic data generation methods~\cite{gao2023magicdrive, zhou2024simgen, hu2024drivingworld, hu2024drivingworld}. 
LDM-based approaches~\cite{lu2024wovogen, swerdlow2024street, wang2024drivedreamer, zhang2024perldiff,  swerdlow2024street, li2025dualdiff} generate diverse driving scene views by conditioning on geometric and textual guidance, ensuring geometric consistency and multi-view alignment.
Recently, subject-driven customization has gained traction in driving scene generation. Inspired by Subject Diffusion~\cite{ma2024subject}, SubjectDrive~\cite{huang2025subjectdrive} introduces a Subject Prompt Adapter, along with a subject bank to enable flexible target replacement for improved sampling diversity.
DriveEditor~\cite{liang2025driveeditor} offers a unified framework for object repositioning, insertion, deletion, and replacement through a learned reconstruction task.
SceneCrafter~\cite{zhu2025scenecrafter} enables the insertion or removal of foreground objects and the modification of global attributes, such as weather and time of day.
MVPbev~\cite{liu2024mvpbev} extends these capabilities with camera-pose editing, supporting generalizable viewpoint changes across different driving camera settings.
However, these methods overlook map editing, limiting the diversity of road topologies presented in synthetic datasets. Our work seeks to fill this gap.

